\section{The BILIP variants: the author's equivalence claim, its analysis, and
consequences for M1}
\label{app:bilip}

This appendix records what is known about the two decoder regularizers that can serve
as BILIP (\cref{sec:pseudoinj:blae}), what the BLAE author said about them, how his
equivalence claim can be read, and what follows for M1. Each statement is one of
three
kinds: \emph{verified} (read in the paper or in the released code), \emph{stated} (by
the author, in correspondence of October 2026), or \emph{reasoning} (mine, from the
theory and the code; not tested, and not claimed by the authors).

\textbf{Note:} In the following $g$ denotes the mixing function of the \emph{head}.

\subsection{What is established about the variants}
\label{app:bilip:facts}
\begin{itemize}[leftmargin=1.5em]
\item \emph{Verified, paper.} The paper defines one bi-Lipschitz regularizer,
Eq.~(12)
      of \cite{zhan2026blae},
      \[
      \mathbb{E}\left[\left(\mathrm{ReLU}\left(\tfrac{1}{\kappa}-\sigma_{\min}\right)\right)^2
      +\left(\mathrm{ReLU}\left(\sigma_{\max}-\kappa\right)\right)^2\right],
      \]
      on the singular values of the decoder Jacobian $J$. Its theory (admissibility,
      embedding dimension $n=O(k)$) is stated for this loss.
\item \emph{Verified, code.} The released code has several \texttt{reg\_grad}
      branches. \texttt{bi-lipschitz\_dec} is Eq.~(12). The default,
      \texttt{bi-lipschitz}, penalizes the off-diagonal entries of $G=J^\top J$ and
      the diagonal entries above $L$, and multiplies its lower-bound term by $0$.
      The encoder branches use a \texttt{tangent\_spaces} attribute that is never
      set, so they cannot run as released.
\item \emph{Verified, notebook.} The released \texttt{swiss\_roll.ipynb}, and an
      earlier version of it, both call \texttt{bi-lipschitz} with \texttt{lam\_grad}
      $=0.3$, \texttt{lam\_geo} $=1$, \texttt{L} $=1$, \texttt{thresh} $=0.3$. These
      are the values of the Swiss Roll column in \cite[Table~3]{zhan2026blae}.
      ($\lambda_{\text{bi-Lip}}=0.3$, $\lambda_{\mathrm{reg}}=1$, $\kappa=1$,
      $\epsilon=0.3$). The BLAE panel of the earlier notebook looks the same, by eye,
      as that of \cite[Figure~3]{zhan2026blae}.
\item \emph{Stated.} The author describes the four branches as follows:
      \texttt{bi-lipschitz} regularizes $J^\top J$ (off-diagonal to zero, diagonal
      capped at $L$); \texttt{bi-lipschitz\_dec} keeps the singular values of the
      decoder Jacobian in $[1/L,L]$; the encoder variant \texttt{bi-lipschitz\_enc} does the same 
      for the encoder Jacobian on the tangent space of the data manifold; 
      the \texttt{bi-lipschitz\_enc+dec} variant constrains only the largest singular value, of 
      both. He states that the paper's experiments use \texttt{bi-lipschitz\_dec}, and that 
      ``theoretically, these four variants are equivalent if 
      $\mathrm{dec}\circ\mathrm{enc}=\mathrm{identity}$ and orthogonal transformations are 
      ignored''; he also claims that empirically \texttt{bi-lipschitz\_dec}
      is used because it works best.
\item \emph{Open.} The notebook and the author's statement point to different
      variants for the Swiss Roll figure. This is not reconciled. It is not known
      which variant produced Figure~3, nor which one the other datasets used.
\end{itemize}

\subsection{The author's equivalence claim}
\label{app:bilip:claim}

\paragraph{Why encoder and decoder constraints coincide (reasoning).}
If $\mathrm{dec}\circ\mathrm{enc}$ is the identity on the data manifold, the chain
rule gives $J_{\mathrm{dec}}J_{\mathrm{enc}}=I$ on the tangent space of the manifold.
On that subspace $J_{\mathrm{dec}}$ is the inverse of $J_{\mathrm{enc}}$, so its
singular values are the reciprocals of those of $J_{\mathrm{enc}}$. The interval
$[1/\kappa,\kappa]$ is mapped to itself by taking reciprocals, hence ``all singular
values of the decoder lie in $[1/\kappa,\kappa]$'' and ``all singular values of the
encoder lie in $[1/\kappa,\kappa]$'' say the same thing. Likewise
$\sigma_{\max}(J_{\mathrm{dec}})\le\kappa$ is the same statement as
$\sigma_{\min}(J_{\mathrm{enc}})\ge1/\kappa$. This is why capping only the maximum,
for
both maps, is enough in the \texttt{enc+dec} variant.

\paragraph{What ``orthogonal transformations are ignored'' means (reasoning).}
Singular values are unchanged by orthogonal maps: $\sigma(QJR)=\sigma(J)$ for
orthogonal
$Q$ (output space) and $R$ (latent space). Being $\kappa$-bi-Lipschitz is also
unchanged, since rotations are isometries. The loss of \cite[Eq.~(12)]{zhan2026blae} therefore 
cannot tell a latent space from the same latent space with rotated axes, and the set of its
minimizers is closed under such rotations. The default variant behaves differently. Under a
latent rotation $G=J^\top J$ becomes $R^\top GR$, so its off-diagonal penalty is zero
only when the latent axes are aligned with the principal axes of the metric.

\paragraph{Why rotations can be ignored for BLAE (reasoning).}
BLAE aims at preserving the geometry of the manifold, and geometry does not depend on
the choice of basis of the latent space: any rotated embedding is as good as any
other. If the rotation of the latent axes may be chosen freely, $G$ can always be
diagonalized, its diagonal is then $\sigma^2$, and the default variant reduces to a
condition on the singular values ($\mathrm{diag}(G)\le L$ means $\sigma_{\max}\le
\sqrt{L}$). In that sense the variants agree.

\paragraph{Where the claim does not carry over (reasoning, from the code).}
\begin{itemize}[leftmargin=1.5em]
\item \emph{The lower bound.} In the code the lower-bound term of \texttt{bi-lipschitz} is 
      multiplied by $0$.
\footnote{Note that there's a commented out line which lower-bounds the diagonal to $1/L$.}
      That variant constrains stretching, not collapse, so it
      is not equivalent to the two-sided loss even up to rotation, and the author's
      description (``stretched or compressed'') does not match it.
\item \emph{Directions the data do not occupy.} Denote $k,\ n$ the intrinsic and the estimated
      latent dimensions, respectively. $\mathrm{dec}\circ\mathrm{enc}=
      \mathrm{id}$ fixes $J_{\mathrm{dec}}$ only on the $k$ directions along the
      encoded manifold. When $n>k$ the other $n-k$ columns are unconstrained by
      reconstruction, while the decoder-side loss still constrains them
      (\cref{app:bilip:m1}).
\footnote{
$\mathrm{dec} \circ \mathrm{enc} = \mathrm{id}$ fixes $J_{\mathrm{dec}}$ only on the $k$ directions 
along the encoded manifold:
\begin{enumerate}[nosep]
\item Let $\Mk$ be the data manifold. Its dimension is $k$, unknown to us but fixed by nature.
\item If $\mathrm{dec}(\mathrm{enc}(x)) = x$ for all $x$ in $\Mk$, differentiate along any curve in
$\Mk$. For every tangent vector $v$ at $x$, $J_{\mathrm{dec}} \cdot J_{\mathrm{enc}} \cdot v = v$.
\item The vectors $J_{\mathrm{enc}} \cdot v$ fill a $k$-dimensional subspace of the latent space $\R^n$, the 
tangent space of the encoded manifold at that point. On that subspace the equation pins down 
$J_{\mathrm{dec}}$. On its orthogonal complement, of dimension $n-k$, the equation says nothing.
\end{enumerate}
}
\item \emph{A fixed basis.} The claim is up to a rotation of the latent axes. M1 does
not
      allow one (\cref{app:bilip:m1}).
\end{itemize}

\subsection{Consequences for M1}
\label{app:bilip:m1}

\paragraph{The basis (reasoning).}
The identifiability result of \cite{ng2025} as applied to M1 promises that style latents are 
recovered up to a coordinate-wise invertible transformation. Each content latent is recovered as a 
function of itself and its surrounding parents.
A coordinate-wise invertible map rescales the columns of the
decoder Jacobian: if $\hat z_i=h_i(z_i)$, the Gram matrices satisfy
$\hat G_{ij}=c_i c_j G_{ij}$ with $c_i=(h_i^{-1})'$ and $G$ the Gram matrix of the
true mixing. A zero off-diagonal entry stays zero and a nonzero one stays nonzero.
So the off-diagonal penalty of \texttt{bi-lipschitz} is met by the true latents only
if
the columns $\partial g/\partial z_i$ are mutually orthogonal, which nothing in M1
guarantees for a nonlinear $g$ and causally dependent content latents. Otherwise the
cheapest way to reduce the penalty is to mix the latents into orthogonalized
combinations, which is what the theory says must not happen, and which competes with
the prior (diagonal Gaussian noise, lower-triangular graph; neither is rotation
invariant). The singular-value loss has no preferred basis, so the prior stays in
charge
of it. This favours \texttt{bi-lipschitz\_dec} for M1.

\paragraph{Unused latent directions, $n>k$ (reasoning).}
The note starts with more content latents than it expects to need
(\cref{sec:schedule:general}), so there will be latent directions that no data vary
along. The loss of \cite[Eq.~(12)]{zhan2026blae} looks at all $n$ singular values of 
$J$ at the encoded points, hence it asks the decoder to respond, with $\sigma\ge1/\kappa$ and
independently of the other columns, to directions the reconstruction does not need. This is
geometrically possible, since $P=8192\gg n$, but it is at odds with the KL term,
which pulls the posterior of an unused latent towards the prior (broad), while a decoder
that reacts to that coordinate perturbs the reconstruction unless the posterior stays
narrow. I expect some tension here; it is not tested. The default variant has no
lower bound and lets the decoder ignore such directions. The paper's embedding statement
concerns a $k$-dimensional manifold in $\R^n$; I do not see that it covers the extra
directions, and I have not checked this further. Ng et al.\ assume that the number of
latents is the true one, so extra latents are outside their theory in any case.

\paragraph{Scale (reasoning).}
Both variants impose an absolute scale on the decoder through $L$ or $\kappa$. The
default $2$ is a guess, it is not on the same scale in the two variants, and it may
conflict with the scale of the IN-VAE latents. The gradient-scale probe
(\cref{sec:schedule:gradlog}) is the tool to check the weight. The weight itself
($10^{-4}$) is far below the paper's ($0.1$ to $0.3$), so at that value the choice of
variant may matter little.

\paragraph{What is implemented.}
Both variants are in \texttt{losses.py}, selected by \texttt{loss.bilip\_variant}
(\cref{sec:schedule:bilipcfg}); the choice is the user's, the configuration default
is
\texttt{bi-lipschitz}. The Eq.~(12) variant takes both extremes from one
\texttt{svdvals} call, where the released code makes two calls that each compute all
singular values.

\subsection{Checks that would settle the open points}
\label{app:bilip:checks}
\begin{itemize}[leftmargin=1.5em]
\item The singular-value spectrum of $J$ at $\hat\mu$ over training. Values pinned at
      $1/\kappa$ point to forced directions, values near $\kappa$ to a scale
      conflict.
\item The posterior variance of each latent, to identify unused latents, and a run
      with a smaller $n_c$ for comparison.
\item An ablation with the BILIP weight set to $0$, on a toy case and on TI, to see
      whether the term does anything. Nothing here shows that it does on the Swiss
      Roll, a flat manifold on which injectivity and reconstruction may already
      suffice.
\item The latent-to-geodesic ratio on held-out pairs, and the stability of the
learned
      graph and of the latents across seeds and under reweighting of the training
      locations, for each variant.
\item The question to the author whether Figure~3 was produced by the released
notebook
      as it stands (\texttt{bi-lipschitz}) or with \texttt{bi-lipschitz\_dec}.
\end{itemize}

\subsection{Status}
\label{app:bilip:status}
Nothing in this appendix has been run on TI data, and the torch code of both variants
is
untested. The statements marked as reasoning are arguments, not results. The account
of
the author's claim is my reconstruction; he did not give a derivation.
